\chapter{Supplementary Question 1 on Power Usage Effectiveness Reporting}
\label{ch:supp:pue}

\section{The question}
\label{sec:supp:pue:q}

\begin{quote}
\itshape
Your submission discusses Singapore's managed capacity allocation regime, where data centre approvals are conditional on meeting a Power Usage Effectiveness (PUE) of 1.25 or better. Do you have data on the typical PUE of data centres currently operating in NSW? Is this information publicly available, and if not, do you support mandatory PUE reporting as a baseline transparency measure?
\end{quote}

In short, Protocol does not hold, and is not aware of, any comprehensive and publicly available dataset reporting the Power Usage Effectiveness of data centres operating in New South Wales. The information is not public because the only fit-for-purpose New South Wales instrument is voluntary and its underlying data are confidential. Protocol supports mandatory reporting as a baseline transparency measure, subject to the qualification that the obligation should capture a basket of measured indicators rather than Power Usage Effectiveness alone. The remainder of this chapter sets out the basis for each of these positions.

\section{What is known about efficiency in operating data centres}
\label{sec:supp:pue:known}

The publicly defensible statements about typical Power Usage Effectiveness are international and national averages, not facility-level New South Wales figures. The most recent global industry survey places the worldwide average at roughly 1.5 and records that it has stagnated rather than improved across recent years \autocite{uptimeGlobalDataCentreSurvey2025}. Commentary on the Australian fleet places the national average near 1.44 for 2023, a figure that itself derives from vendor analysis rather than any official register \autocite{nlyteAustralianDataCentreEfficiency2024}. These averages conceal a wide distribution. The most efficient hyperscale operators report fleet-wide values close to 1.1, and they can do so credibly only because they measure continuously, explaining that \enquote{Our calculations are based on continuously measuring the entire worldwide fleet performance of our data centers throughout the year} \autocite{googleDataCentreEfficiency2025}. Singapore's condition of a Power Usage Effectiveness of 1.25 or better therefore sits below the Australian average and near the upper boundary of best practice, which is to say it is demanding but demonstrably achievable.

These figures should be read against the trajectory of demand that makes efficiency salient. Artificial intelligence has ended the long plateau in which efficiency gains offset workload growth, and data centres \enquote{have become one of the fastest-growing electricity consumers globally} \autocite{shengPowerAIData2026}, with the underlying compute exhibiting \enquote{a Deep Learning Era compute doubling time of around 6 months} \autocite{sevillaComputeTrendsThree2022}. The resulting consumption is also systematically under-stated by simple proxies, since \enquote{naive energy estimates based on FLOPs or theoretical GPU utilization significantly underestimate real-world energy consumption} \autocite{fernandezEnergyConsiderationsLarge2025a}. Reliable measurement at the facility, rather than estimation from first principles, is therefore the only sound basis for policy.

\section{Why no New South Wales figure is available}
\label{sec:supp:pue:gap}

The absence of a New South Wales figure is structural rather than incidental. New South Wales has access to a purpose-built instrument, the NABERS Energy for Data Centres scheme, which rates facilities from one to six stars on measured operational performance across three rating types, being Infrastructure, Information Technology Equipment, and Whole Facility \autocite{nabersDataCentres2024}. Participation in that scheme is voluntary. The scheme maintains a public register of certified ratings, but its coverage is partial because only operators who choose to participate and to publish appear in it, and the underlying energy data remain confidential \autocite{nabersEnergyDataCentresRules2024}. The aggregate efficiency of the New South Wales fleet is consequently unobservable.

The mandatory reporting that does exist is narrow. From 1 July 2025 the Commonwealth requires a minimum five-star NABERS Energy rating for data centres hosting Australian Government workloads, an obligation that binds only that subset of facilities \autocite{lsjDataCentresFiveStar2024}. The National Greenhouse and Energy Reporting scheme captures total energy and emissions above a facility threshold of 100 terajoules but does not measure or require Power Usage Effectiveness, because it is an emissions accounting framework rather than an efficiency standard \autocite{dcceewNGERScheme2025}. Proposed expansion of the Commercial Building Disclosure programme to data centres is the most likely future vehicle, advanced on the basis that \enquote{The commercial building sector is responsible for around 25\% of electricity use and 10\% of total emissions in Australia}, but it remains in a consultation and roadmap phase and offers no near-term remedy \autocite{dcceewCBDExpansion2024}. The net position is that private operators outside the Commonwealth perimeter face no obligation to disclose efficiency at all.

\section{International precedent for mandatory disclosure}
\label{sec:supp:pue:precedent}

The feasibility of mandatory disclosure is settled by precedent. The European Union has moved from voluntary benchmarking to a statutory regime requiring \enquote{monitoring and reporting of the energy performance of data centres}, implemented through a \enquote{Delegated Regulation on the establishment of a common Union rating scheme for data centres (EU/2024/1364)} \autocite{ecEnergyPerformanceDataCentres2024}. That regime is explicitly a basket rather than a single ratio, since the directive \enquote{requires data centers to report operational efficiency metrics such as power usage effectiveness (PUE) and water usage effectiveness (WUE), and adopt measures that optimize electricity and water usage} \autocite{datacenterKnowledgeCompliance2026}. The contrast with the United States is instructive, where \enquote{In the US, no federal equivalent of the EED exists, but state-level activity is picking up} \autocite{datacenterKnowledgeCompliance2026}. The emerging state instruments couple disclosure to network management, as with the Oregon measure that \enquote{establishes a special electricity rate for data centers and other large power consumers, incentivizing efficiency and grid-friendly load profiles} \autocite{datacenterKnowledgeCompliance2026}.

There is an instructive parallel in the governance of computing more generally, where reporting is treated as the precondition for any further lever. Compute is considered \enquote{a particularly effective point of intervention} precisely because \enquote{it is detectable, excludable, and quantifiable} \autocite{sastryComputingPowerGovernance2024}, and the frontier-regulation literature places \enquote{registration and reporting requirements to provide regulators with visibility into frontier AI development processes} ahead of substantive obligations \autocite{anderljungFrontierAIRegulation2023}. The same logic recommends that an efficiency standard for data centres begin with measurement and disclosure, not with an unenforceable target.

\section{Why disclosure is necessary but not sufficient}
\label{sec:supp:pue:nec}

Two qualifications discipline the recommendation. First, disclosure is necessary for accountability but not sufficient for outcomes, and it is useful to separate the two functions that transparency performs. The distinction in the governance literature \enquote{between principles of \enquote{fishbowl transparency} and \enquote{reasoned transparency}} captures the point, since publishing raw ratings serves public scrutiny while reporting to the regulator serves enforcement, and a well-designed regime needs both \autocite{coglianeseTransparencyAlgorithmicGovernance2019a}. An efficiency-conditioned regime in the manner of Singapore does not strictly require a public register in order to function, because measurement reported to and verified by the regulator would suffice for enforcement; public disclosure is justified separately, as the layer that allows the standard and its administration to be scrutinised.

Second, Power Usage Effectiveness alone is the wrong target. It is gameable, since a facility can improve the ratio by raising its information technology baseline, and it is blind to absolute consumption, carbon and water. The literature accordingly advocates \enquote{transparency mechanisms, such as disclosing the GHG footprint of AI systems}, situating efficiency within a broader resource and emissions frame in which \enquote{The ICT sector contributes up to 3.9 percent of global greenhouse gas (GHG) emissions-more than global air travel at 2.5 percent} \autocite{hackerSustainableAIRegulation2024}. The case for disclosing emissions specifically is strengthened by evidence that estimates omitted from primary reporting are wildly unreliable, with one analysis finding that \enquote{Estimates of emissions in papers that omitted them have been off 100x-100,000x} \autocite{pattersonCarbonFootprintMachine2022a}. Water reinforces the same conclusion, since the most consequential figures have historically been concealed, as when \enquote{training the GPT-3 language model in Microsoft's state-of-the-art U.S. data centers can directly evaporate 700,000 liters of clean freshwater, but such information has been kept a secret} \autocite{liMakingAILess2025}. A mandate confined to Power Usage Effectiveness would license precisely the partial visibility that has allowed these costs to go unmeasured.

\section{Recommendation}
\label{sec:supp:pue:rec}

Protocol recommends that New South Wales make disclosure a condition of development consent and of continued operation for data centres above a defined information technology power threshold, calibrated on the European Union's 500 kilowatt model. The disclosed basket should comprise the NABERS Whole Facility rating and its underlying Power Usage Effectiveness, water usage effectiveness, absolute annual energy, and renewable share, reported annually to the regulator and published in aggregate and at facility level. NABERS already supplies the administrative spine and the measurement methodology, so the missing element is solely the obligation to measure and disclose. This sequencing is the precondition for every downstream lever the Committee may later consider, including a Singapore-style allocation of capacity against an efficiency standard, because capacity cannot be allocated against a benchmark that is neither measured nor published, and a target that cannot be audited cannot be enforced \autocite{imanHiddenCostsIntelligence2025, hackerSustainableAIRegulation2024}.
